\section*{Chapitre \ref{ch:mx:build-an-agent-based-market} --- Construction~: un marché multi-agents}

\begin{solution}{exo:mx:build-an-agent-based-market:1}
$(0.212-0.158)/0.046=1.17$ écarts types, ce qui ajoute $1.17^2=1.4$ à la distance (un quart de la distance inédite de 5.4).
\end{solution}

\begin{solution}{exo:mx:build-an-agent-based-market:2}
Comme $\tau^{-(\alpha-1)}$~: $\tau^{-0.5}$ pour une queue de 1.5 et $\tau^{-1.5}$ pour 2.5. Seule la première est à mémoire longue ($\alpha<2$~: les autocorrélations
ne sont pas sommables)~; la queue calibrée de 2.5 donne une mémoire courte.
\end{solution}

\begin{solution}{exo:mx:build-an-agent-based-market:3}
L'écart normalisé de chaque fait a pour variance $1/6+1/8=0.29$, donc la distance espérée vaut $7\times0.29=2.0$. Le plancher, calculé sur 6 heures inédites de la cible face à la
moyenne de 8 heures, vaut 1.7, proche de la même espérance~; la population calibrée atteint 5.4, soit environ 2.7 fois ce niveau~:
un désajustement au-delà du bruit d'échantillonnage, surtout sur le regroupement de volatilité, la profondeur et la mémoire de signe à retard 1.
\end{solution}

\begin{solution}{exo:mx:build-an-agent-based-market:4}
Deux raisons. La moyenne sur trois heures porte plus de bruit de simulation que celle sur six ($7\times(1/3+1/8)=3.2$ attendu contre 2.0), et le point a été
retenu parce que ses trois heures sont tombées les plus proches parmi 27, si bien que sa distance dans l'échantillon reflète aussi quelles heures elles étaient~: sur les
heures de calibration, son kurtosis valait 3.11, contre 1.90 sur les heures inédites. La distance sur les heures inédites est la seule honnête.
\end{solution}

\begin{solution}{exo:mx:build-an-agent-based-market:5}
Sans eux, rien ne fait bouger le cours entre les ordres de bruit, que le carnet profond des fournisseurs de liquidité absorbe~: le prix milieu reste immobile pendant la plupart des
intervalles de 30 secondes, et ne bouge que lorsqu'une série d'ordres de bruit ou une rafale de chartistes vide un niveau. Des rendements presque toujours nuls, avec de rares sauts, ont un kurtosis très
élevé. Ce sont les fondamentalistes qui font suivre $V$ au cours par petits pas.
\end{solution}

\begin{solution}{exo:mx:build-an-agent-based-market:6}
Ses exécutions subissent une sélection adverse~: les fondamentalistes traitent contre lui quand $V$ a bougé, si bien que, sans skew, son inventaire croît dans le mauvais sens et
y reste tant que le cours continue de s'éloigner de lui. Le skew fait que ses cotations attirent les transactions qui réduisent l'inventaire, de sorte qu'il détient moins (1.6 lots en moyenne quadratique pour un skew de 0.3, contre 14.5), moins longtemps, et empoche la fourchette sur un flux qui se renouvelle.
\end{solution}

\begin{solution}{exo:mx:build-an-agent-based-market:7}
Oui. Sur les cinq faits, le meilleur point de la grille devient un taux de bruit de 1.5, une queue des séries de 2.5 et un taux de chartistes de 0.1, à une distance de 5.1 sur les
heures de calibration~; le point du chapitre arrive en deuxième position avec 7.7. Le flux de bruit plus rapide amincit le carnet (profondeur d'environ 9 lots contre 24 pour la cible) et élargit
la fourchette (environ 1.08 ticks contre 1.04, soit plus de deux écarts types de la cible)~: la calibration complète l'a rejeté à cause de son carnet, que les cinq
faits ne voient plus. Le choix des faits décide de la réponse.
\end{solution}

\begin{solution}{exo:mx:build-an-agent-based-market:8}
Ajuster des faits n'identifie pas les mécanismes~: des populations différentes peuvent produire les mêmes sept moments (ici, retirer les chartistes n'a modifié aucun fait
de façon significative, si bien que toute part de chartistes proche de la part calibrée s'ajuste tout aussi bien), les faits non ajustés montrent que la population manque de quelque chose que la cible
possède, et la cible est elle-même un modèle. La calibration dit que la population est suffisante pour certaines questions, non ce que sont réellement les traders.
\end{solution}

\begin{solution}{pb:mx:build-an-agent-based-market:1}
\textbf{1.} Voir les définitions.
\textbf{2.} Lux et Marchesi~: les interactions entre fondamentalistes et chartistes produisent des queues épaisses et une volatilité regroupée à partir de nouvelles qui n'en ont ni l'une ni l'autre.
Farmer, Patelli et Zovko~: un placement aléatoire des ordres explique 96~\% de la variance transversale des fourchettes et 76~\% de celle de la diffusion des cours sur 11
actions cotées à Londres avec un seul paramètre libre.
\textbf{3.} Des \omterm{def:mx:the-limit-order-book:orders}{ordres à cours limité} d'un lot de chaque côté, à un demi-tick du meilleur cours opposé et le reste jusqu'à dix ticks derrière celui-ci, chacun annulé à un
taux constant~; des \omterm{def:mx:the-limit-order-book:orders}{ordres au marché} à taux constant avec des signes en séries de Pareto~; des \omterm{def:mx:the-limit-order-book:orders}{ordres au marché} vers la valeur cachée à un taux proportionnel à l'écart~;
des \omterm{def:mx:the-limit-order-book:orders}{ordres au marché} suivant la tendance sur 30 secondes, à un taux proportionnel à sa taille.
\textbf{4.} Eux seuls voient $V$.
\textbf{5.} Pour que des simulations qui ne diffèrent que d'un agent partagent les mêmes tirages (nombres aléatoires communs).
\textbf{6.} Voir la définition~; deux lots aux meilleurs cours des autres (à un tick à l'intérieur lorsque leur fourchette dépasse deux ticks), décalés du skew multiplié par
l'inventaire, et aucun cours acheteur (vendeur) au-delà de 20 lots en position longue (courte).
\textbf{7.} Chaque nouveau cours acheteur était fixé sous son propre cours acheteur précédent en position longue, si bien que le cours acheteur glissait vers le bas à chaque rafraîchissement.
\textbf{8.} 14.5, 4.2 et 1.6 lots~; $-116$, 448 et 545 dollars de l'heure pour des skews de 0, 0.1 et 0.3.
\textbf{9.} De grands ordres à loi de puissance, découpés en enfants à taux constant, donnent une autocorrélation de signe qui décroît comme $\tau^{-(\alpha-1)}$~; les séries de signes des preneurs de liquidité de bruit
ont des longueurs de Pareto de queue $\alpha$.
\textbf{10.} Un excès dû aux séries, qui décroît avec une pente de $-1.15$ pour les retards 1 à 7 (la loi prévoit $-1.5$), disparaît au retard 10~; le reste vient des
fondamentalistes.
\textbf{11.} Voir la définition~; pour que les écarts entre points de la grille ne soient pas des écarts de hasard.
\textbf{12.} Huit heures de \texttt{firm.tape} sans nouvelles~; kurtosis et autocorrélation à retard 1 des rendements sur 30 secondes, regroupement de volatilité, autocorrélation de signe aux
retards 1 et 10, fourchette, profondeur.
\textbf{13.} Bruit 0.9, 1.2, 1.5~; queue des séries aucune, 1.5, 2.5~; chartistes 0, 0.05, 0.1~; meilleur 1.2, 2.5, 0.05.
\textbf{14.} \emph{Résultat nommé}~: après calibration, la distance vaut 10.1 sur les trois heures de calibration et 5.4 sur six heures inédites, contre 27.9 pour
le point de départ et un plancher de 1.7 pour le marché cible lui-même~; les séries de signes portent la mémoire de signe à retard 1 ($-2.77$, erreur type 0.15), les fondamentalistes le
mouvement du cours (kurtosis $+8.6$ sans eux), le teneur de marché la fourchette ($+9.4$), et les chartistes rien de significatif ($+1.11$, avec une erreur type de 0.55).
\textbf{15.} Regroupement de volatilité, profondeur et mémoire de signe à retard 1~: la cible a un niveau d'activité aléatoire et un flux regroupé à la Hawkes que la population n'a pas.
\textbf{16.} Les configurations d'agents interactives étaient plus réalistes que le rejeu, et davantage encore avec une valeur fondamentale tirée de données historiques.
\textbf{17.} Un tape rejoué ne réagit pas aux ordres de l'algorithme, si bien que l'impact et la réaction des autres manquent.
\textbf{18.} Plusieurs populations ajustent les mêmes faits~; un type d'agent peut être retiré sans grand changement.
\textbf{19.} Il faut s'en servir de façon comparative, algorithme contre algorithme avec les mêmes graines, et vérifier que les conclusions ne dépendent pas des faits qu'il omet.
\textbf{20.} Pour répondre aux questions «~que se passerait-il si~» sur les mécanismes d'un marché qui réagit.
\end{solution}

\begin{solution}{iq:mx:build-an-agent-based-market:1}
Des agents qui placent et annulent des \omterm{def:mx:the-limit-order-book:orders}{ordres à cours limité} et au marché à des taux aléatoires, sans stratégie~; avec les taux de flux d'ordres en entrée, il prédit bien les fourchettes et
la diffusion des cours d'une action à l'autre (Farmer, Patelli et Zovko), ce qui montre quelle part le seul mécanisme explique.
\end{solution}

\begin{solution}{iq:mx:build-an-agent-based-market:2}
Choisir les faits qui comptent pour la question (rendements, flux d'ordres, carnet), les mesurer sur l'action avec leur variabilité d'échantillonnage, et minimiser une
distance pondérée sur les paramètres de la population par la \omterm{def:mx:build-an-agent-based-market:msm}{méthode des moments simulés}, avec des nombres aléatoires communs~; vérifier l'ajustement sur des journées simulées inédites et sur
des faits non utilisés pour l'ajustement, et comparer la distance à ce que le seul bruit d'échantillonnage donne.
\end{solution}

\begin{solution}{iq:mx:build-an-agent-based-market:3}
Une seule file d'événements ordonnée par le temps et une séquence de départage, chaque flux aléatoire initialisé par agent et par usage, aucune horloge murale ni temps de threads dans
la logique~; ainsi, les deux algorithmes font face au même flux exogène et leur écart n'est pas du bruit.
\end{solution}

\begin{solution}{iq:mx:build-an-agent-based-market:4}
Pour rendre plus probables les exécutions qui réduisent l'inventaire et moins probables celles qui l'alourdissent, ce qui réduit le risque d'inventaire et la sélection adverse~; assez pour que
l'inventaire revienne vers sa moyenne dans le délai de détention que le desk accepte, ce qui, sur le marché du chapitre, représentait 0.1 à 0.3 tick par lot.
\end{solution}

\begin{solution}{iq:mx:build-an-agent-based-market:5}
Il apprend les particularités du simulateur~: des agents qui ne s'y adaptent pas, une liquidité qui se reconstitue selon des règles fixes, des faits absents comme le regroupement de volatilité~; il peut
exploiter des régularités absentes de la réalité. Randomiser la population, tester sur des configurations de test et sur des données réelles avant de négocier.
\end{solution}

\begin{solution}{iq:mx:build-an-agent-based-market:6}
Quels faits le simulateur ajuste et lesquels il manque, si l'avantage de l'algorithme tient face à d'autres populations plausibles et à d'autres graines, si son comportement
dans des scénarios de stress (suspensions de cotation, sauts de cours, flux à sens unique) a été testé, et un plan pour un essai réel limité, avec des contrôles.
\end{solution}
